% Options for packages loaded elsewhere
\PassOptionsToPackage{unicode}{hyperref}
\PassOptionsToPackage{hyphens}{url}
\PassOptionsToPackage{dvipsnames,svgnames,x11names}{xcolor}
%
\documentclass[
  usletterpaper,
]{article}
\usepackage{amsmath,amssymb}
\usepackage{setspace}
\usepackage{iftex}
\ifPDFTeX
  \usepackage[T1]{fontenc}
  \usepackage[utf8]{inputenc}
  \usepackage{textcomp} % provide euro and other symbols
\else % if luatex or xetex
  \usepackage{unicode-math} % this also loads fontspec
  \defaultfontfeatures{Scale=MatchLowercase}
  \defaultfontfeatures[\rmfamily]{Ligatures=TeX,Scale=1}
\fi
\usepackage{lmodern}
\ifPDFTeX\else
  % xetex/luatex font selection
  \setmainfont[]{Libertinus Serif}
  \setsansfont[]{Libertinus Sans}
  \setmonofont[]{Iosevka}
  \setmathfont[]{Libertinus Math}
\fi
% Use upquote if available, for straight quotes in verbatim environments
\IfFileExists{upquote.sty}{\usepackage{upquote}}{}
\IfFileExists{microtype.sty}{% use microtype if available
  \usepackage[]{microtype}
  \UseMicrotypeSet[protrusion]{basicmath} % disable protrusion for tt fonts
}{}
\makeatletter
\@ifundefined{KOMAClassName}{% if non-KOMA class
  \IfFileExists{parskip.sty}{%
    \usepackage{parskip}
  }{% else
    \setlength{\parindent}{0pt}
    \setlength{\parskip}{6pt plus 2pt minus 1pt}}
}{% if KOMA class
  \KOMAoptions{parskip=half}}
\makeatother
\usepackage{xcolor}
\usepackage[margin=1.3in]{geometry}
\usepackage{longtable,booktabs,array}
\usepackage{calc} % for calculating minipage widths
% Correct order of tables after \paragraph or \subparagraph
\usepackage{etoolbox}
\makeatletter
\patchcmd\longtable{\par}{\if@noskipsec\mbox{}\fi\par}{}{}
\makeatother
% Allow footnotes in longtable head/foot
\IfFileExists{footnotehyper.sty}{\usepackage{footnotehyper}}{\usepackage{footnote}}
\makesavenoteenv{longtable}
\setlength{\emergencystretch}{3em} % prevent overfull lines
\providecommand{\tightlist}{%
  \setlength{\itemsep}{0pt}\setlength{\parskip}{0pt}}
\setcounter{secnumdepth}{5}
\ifLuaTeX
\usepackage[bidi=basic]{babel}
\else
\usepackage[bidi=default]{babel}
\fi
\babelprovide[main,import]{english}
\ifPDFTeX
\else
\babelfont{rm}[]{Libertinus Serif}
\fi
% get rid of language-specific shorthands (see #6817):
\let\LanguageShortHands\languageshorthands
\def\languageshorthands#1{}
\usepackage[cache=false]{minted}
\setminted{frame=single, framesep=2mm, linenos, breaklines}
\usepackage[font={footnotesize}, margin=1cm, labelfont=bf]{caption}
\usepackage{subcaption}
\usepackage{multicol}
\usepackage{pgfplots}
\pgfplotsset{compat=newest, width=10cm}
\usepgfplotslibrary{groupplots}
\usepackage{extarrows}
\usepackage[sf,bf]{titlesec}
\usepackage{tikz}
\makeatletter
\let\listoflistings\@undefined
\makeatother
\makeatletter
\@ifpackageloaded{subfig}{}{\usepackage{subfig}}
\@ifpackageloaded{caption}{}{\usepackage{caption}}
\captionsetup[subfloat]{margin=0.5em}
\AtBeginDocument{%
\renewcommand*\figurename{Figure}
\renewcommand*\tablename{Table}
}
\AtBeginDocument{%
\renewcommand*\listfigurename{List of Figures}
\renewcommand*\listtablename{List of Tables}
}
\newcounter{pandoccrossref@subfigures@footnote@counter}
\newenvironment{pandoccrossrefsubfigures}{%
\setcounter{pandoccrossref@subfigures@footnote@counter}{0}
\begin{figure}\centering%
\gdef\global@pandoccrossref@subfigures@footnotes{}%
\DeclareRobustCommand{\footnote}[1]{\footnotemark%
\stepcounter{pandoccrossref@subfigures@footnote@counter}%
\ifx\global@pandoccrossref@subfigures@footnotes\empty%
\gdef\global@pandoccrossref@subfigures@footnotes{{##1}}%
\else%
\g@addto@macro\global@pandoccrossref@subfigures@footnotes{, {##1}}%
\fi}}%
{\end{figure}%
\addtocounter{footnote}{-\value{pandoccrossref@subfigures@footnote@counter}}
\@for\f:=\global@pandoccrossref@subfigures@footnotes\do{\stepcounter{footnote}\footnotetext{\f}}%
\gdef\global@pandoccrossref@subfigures@footnotes{}}
\@ifpackageloaded{float}{}{\usepackage{float}}
\floatstyle{ruled}
\@ifundefined{c@chapter}{\newfloat{codelisting}{h}{lop}}{\newfloat{codelisting}{h}{lop}[chapter]}
\floatname{codelisting}{Listing}
\newcommand*\listoflistings{\listof{codelisting}{List of Listings}}
\makeatother
\ifLuaTeX
  \usepackage{selnolig}  % disable illegal ligatures
\fi
\usepackage[style=ieee,]{biblatex}
\addbibresource{bibliography.bib}
\nocite{*}
\IfFileExists{bookmark.sty}{\usepackage{bookmark}}{\usepackage{hyperref}}
\IfFileExists{xurl.sty}{\usepackage{xurl}}{} % add URL line breaks if available
\urlstyle{tt}
\hypersetup{
  pdftitle={CPPDescent: A C++ library for the creation of K-NN graphs from multi-dimensional datasets},
  pdfauthor={Konstantinos Chousos; Anastasios-Phaedon Seitanidis},
  pdflang={en},
  colorlinks=true,
  linkcolor={Maroon},
  filecolor={cyan},
  citecolor={red},
  urlcolor={blue},
  pdfcreator={LaTeX via pandoc}}

\title{CPPDescent: A C++ library for the creation of K-NN graphs from
multi-dimensional datasets}
\usepackage{etoolbox}
\makeatletter
\providecommand{\subtitle}[1]{% add subtitle to \maketitle
  \apptocmd{\@title}{\par {\large #1 \par}}{}{}
}
\makeatother
\subtitle{Software Development for Computing Systems, Winter
2023-2024\\\medskip Department of Informatics and Telecommunications, University of Athens}
\author{Konstantinos
Chousos\thanks{ID:~1115202000215} \and Anastasios-Phaedon
Seitanidis\thanks{ID:~1115202000179}}
\date{\today}

\begin{document}
\maketitle

\setstretch{1.1}
\begin{abstract}{}{}%
This library is a C++ implementation of the ``NN-Descent'' algorithm by
Dong et al. \autocite{dongEfficientKnearestNeighbor2011}, with a few
improvements and optimisations (e.g.~usage of random projection trees
\autocite{dasguptaRandomProjectionTrees2008}). This implementation
serves as an entry for the
\href{https://2023.sigmod.org/sigmod_student_research_competition.shtml}{ACM
SIGMOD 2023 competition}. It was developed as part of the \emph{Software
Development for Computing Systems} course of the Department of
Informatics and Telecommunications, taught in the winter of 2023
\autocite{ioannidisAnaptyxiLogismikoyGia}. The source code is available
at \url{https://github.com/kchousos/CPPDescent}.%
\end{abstract}

\begin{multicols}{2}

\section{Introduction}\label{introduction}

As is tradition for this course, our project for the semester was the
challenge of last spring's ACM SIGMOD Student Research Competition. The
theme of the competition was the efficient and fast creation of a
\emph{K Nearest Neighbor} graph, where its vertices belong to a dataset
of 100-dimensions vectors/datapoints of floats.

A brute force approach to this problem has a complexity of \(O(n^2)\),
since each datapoint needs to check all others so that it can \emph{weed
out} the K points nearest to it. An answer to this problem was given in
2011 by Dong et al.~in their paper titled ``Efficient K-Nearest Neighbor
Graph Construction for Generic Similarity Measures''
\autocite{dongEfficientKnearestNeighbor2011}, where they present the
\emph{NN-Descent} algorithm. This algorithm is based at the very simple
idea that \emph{``a neighbor of a neighbor is also likely to be a
neighbor''}. According to them, this algorithm is shown to have a
complexity of \(O(n^{1.14})\). This is a major step-up from the
quadratic time, especially when it comes to big dataset.

Our library does not only implement the NN-Descent algorithm. It also
implements other optimizations that e.g.~make the starting graph more
like the one we want to end up with, lowering this way the number of
iterations the algorithm must do to reach a satisfying result. These
optimizations are discussed in sec.~\ref{sec:optimizations}.

\subsection{Project layout}\label{project-layout}

\begin{itemize}
\tightlist
\item
  The library itself is stored in the \mintinline[]{text}{src/}
  directory. It is compiled to a \mintinline[]{text}{.a} file, so that
  it can be linked with any executable.
\item
  The main app that results to the main executable is in the
  \mintinline[]{text}{app/} directory. This file contains the only
  \mintinline[]{text}{main()} on the project, and can be thought as the
  API between the user and the library itself.
\item
  The directory \mintinline[]{text}{datasets/} contains a lot of the
  datasets of the SIGMOD competition in binary format. The subdirectory
  \mintinline[]{text}{computed/} contains some of the K-NN graphs that
  were computed using the brute force algorithm for comparison purposes.
  They are also in binary format. For which dataset and for which K they
  correspond to can be deduced from the respective filename.
\item
  The \mintinline[]{text}{docs/} subdirectory contains the
  \mintinline[]{text}{html/} and \mintinline[]{text}{latex/}
  subdirectories that they themselves contain files respective to each
  format. It is documentation of the code, automatically generated by
  code comments using \href{https://www.doxygen.nl/}{Doxygen}. The
  \mintinline[]{text}{lcov/} subdirectory contains files that are
  generated by \href{https://github.com/linux-test-project/lcov}{LCOV}
  and are used to monitor the code coverage provided by the tests (see
  below)\footnote{At the moment, the code coverage of the
    \mintinline[]{text}{cppdescent} library itself is pretty low, since
    it went very recently under a large refactoring and overhaul.},
\item
  The \mintinline[]{text}{extern/} directory contains the source code
  for the \href{https://github.com/google/googletest}{Google Test}
  suite. It is included as a git submodule.
\item
  The \mintinline[]{text}{include/} directory contains header files for
  the library and the home-made data structures it uses.
\item
  The \mintinline[]{text}{test/} directory contains unit tests for the
  library.
\item
  The \mintinline[]{text}{helper.sh} file is a bash script that is used
  to automate some repetitive tasks, such as execution of all the tests,
  or documentation generation.
\item
  The \mintinline[]{text}{experiments.sh} is a script that executes the
  main program in quiet mode, and outputs a \mintinline[]{text}{.csv}
  file that contains the results of different executions for different
  parameters (see sec.~\ref{sec:experiments}).
\end{itemize}

\subsection{Data structures}\label{data-structures}

\section{Optimizations}\label{sec:optimizations}

\section{Experiments}\label{sec:experiments}

\subsection{Experimental Setup}\label{experimental-setup}

\subsubsection{Performance Measures}\label{performance-measures}

\subsubsection{Default Parameters}\label{default-parameters}

\subsubsection{System Environment}\label{system-environment}

\subsection{Experimental Results}\label{experimental-results}

\subsubsection{Performance}\label{performance}

\paragraph{Execution Speed}\label{execution-speed}

\paragraph{Memory Consumption}\label{memory-consumption}

\paragraph{Single vs.~Multi-Threading}\label{single-vs.-multi-threading}

\subsubsection{Performance as Data
Scales}\label{performance-as-data-scales}

\subsubsection{Tweaking the Parameters}\label{tweaking-the-parameters}

\section{Conclusion}\label{conclusion}

\section{Acknowledgements}\label{acknowledgements}

\end{multicols}

\printbibliography

\end{document}
